% !TeX root=../main.tex
\chapter{جمع‌بندی و کارهای پیش‌رو}
%\thispagestyle{empty} 
\section{مسیر پژوهش و دستاورد‌‌ها}
این پژوهش با الهام از ساختار و کارایی شگفت‌انگیز حس لامسه انسان، یک مسیر تحقیقاتی از طراحی و ساخت نمونه‌های اولیه تا توسعه و ارزیابی یک سیستم حسگر لامسه‌ِ چندوجهی یکپارچه و هوشمند برای کاربردهای رباتیک طی نمود. هدف غایی، فراتر رفتن از محدودیت‌های حسگرهای تک‌وجهی و دستیابی به درکی غنی و قابل اعتماد از تعاملات فیزیکی ربات با اشیاء، به ویژه در وظایف ظریف و پیچیده‌ای مانند گرفتن و دستکاری اشیاء نرم بود.

\subsection{حسگر مبتنی بر مکانیزم بارومتریک}

مسیر پژوهش با طراحی و ساخت یک حسگر لامسه تک‌‌محوره مبتنی بر مبدل فشار بارومتریک مدل \lr{MPL3115A2} آغاز شد. این انتخاب به دلیل مزایای کلیدی این مبدل‌ها از جمله دقت ذاتی بالا، ابعاد کوچک، هزینه پایین و سادگی در پیاده‌سازی الکترونیکی صورت گرفت. فرآیند ساخت، شامل یک نوآوری کلیدی در اصلاح مکانیکی محفظه فلزی حسگر بود که با افزایش کوپلینگ مکانیکی بین لایه کشسان سیلیکونی و دیافراگم داخلی، حساسیت و خطی بودن پاسخ حسگر را به شکل چشمگیری بهبود بخشید. مجموعه کالیبراسیون ساده اما دقیقی برای ارزیابی عملکرد این نمونه اولیه پیاده‌سازی شد و نتایج، عملکردی بسیار مطلوب با مشخصه خطی بالا، پسماند بسیار ناچیز و دقت قابل توجه را به اثبات رساند. صحت اندازه‌گیری نیرو در بازه کاری ۰ تا ۵.۶ نیوتن به کمتر از $1.54\%$ از کل بازه و دقت آن در حدود $2.8mN$ ارزیابی شد. موفقیت این بخش از پژوهش و نوآوری‌های ارائه شده در آن، منجر به پذیرش و انتشار مقاله‌ای
\cite{sheykh2025Design}
 در سی و سومین کنفرانس بین‌المللی مهندسی برق ایران (ICEE) گردید. در ادامه این مرحله، تلاش شد تا با استفاده از آرایه‌ای متشکل از سه حسگر بارومتریک، قابلیت اندازه‌گیری نیروی سه‌محوره محقق شود. با این حال، این رویکرد به دلیل آسیب‌پذیری بسیار بالای قطعات در حین فرآیند حساس اصلاح مکانیکی و اشباع دائمی آن‌ها پس از این فرآیند، عملاً با شکست مواجه شد و به عنوان یک تجربه، توسعه آن به کارهای آتی با مبدل‌های مناسب‌‌تر موکول گردید.

\subsection{حسگر یکپارچه‌ی چندوجهی}
با الهام از ساختار لامسه‌ای انگشتان انسان و استناد به پژوهش‌های معرفی شده در \ref{sec:mechanisms} روش‌های تبدیل زیر نامزد تجمیع در یک برد مدار چاپی یکپارچه شدند. مزیت بنیادین این رویکرد، غلبه بر محدودیت‌های ذاتی هر مکانیزم به تنهایی و دستیابی به یک جریان داده غنی و جامع بوده که تصویری کامل‌تر از رویداد تماس فیزیکی ارائه می‌دهد. علی‌رغم ترکیب داده‌ی حسگرهای انتخاب شده بر بستر یک مدل شبکه عصبی، هر کدام از این حسگر‌ها با هدفی انتخاب شدند که به شرح زیر است:
\begin{itemize}
	\item 
	\textbf{حسگر بارومتریک \lr{MPL3115A2} :}
	این حسگر به عنوان نماینده‌ی مکانیزم تبدیل پیزومقاومتی معرفی شده در\ref{sec:mechanisms:res} انتخاب شد؛ حسن عملکرد این حسگر در اندازه‌گیری نیروی عمودی با دقت بالا در فصل \ref{chap2} نشان داده‌شد. از آنجایی که در سناریو‌های رباتیک و به خصوص گرفتن اشیاء، سنجش نیروی عمودی اهمیت بسزایی دارد، این قسمت از حسگر لامسه با این ایده در نظر گرفته شد تا بتواند سنجش دقیقی از این مولفه‌ی نیرو داشته باشد.
	\item 
	\textbf{حسگر اثر هال \lr{MMC5883MA}:}
	این حسگر یک سنجشگر میدان مغناطیسی در سه محور است، با الهام از پژوهش‌های معرفی شده در \ref{sec:mechanisms:mag} و بهره بردن از یک آهنربای محاط شده در محفظه‌ی کشسان بالای قطعه، این حسگر جابجایی های آهنربا که متاثر از اعمال نیرو به سطح کشسان بوده را به صورت تغییر در میدان مغناطیسی منعکس می‌کرد. در نتیجه‌، حساسیت این حسگر به نیرو‌های اعمال شده در جهات مختلف دلیل انتخاب آن به عنوان یکی دیگر از مکانیزم‌های تبدیل در حسگر لامسه بود. از آنجایی که نیروی مماسی ناشی از وزن اشیاء نیز اهمیت بالایی دارد. مطابق شکل 
	\ref{fig:components3d}
	، این حسگر در زیر یکی از قطب‌های آهنربای رینگی شعاعی نصب شد، به‌طوری ‌که خطوط میدان آهنربا تقریباً عمود بر حسگر باشد. این جاگذاری باعث شد که تغییرات میدان در محور Z حسگر به دلیل جابجایی اندک آهنربا در اثر نیروی مماسی  بیشینه شود. بنابراین، \lr{MMC5883MA} حساسیت بیشتری نسیت به نیرو های مماسی داشته و نقش کلیدی در ثبت مولفه‌های برشی نیرو ایفا می‌کند.
	 \item
	\textbf{حسگر اثر هال \lr{QMC5883L}:}
	 عملکرد و وجه انتخاب این حسگر مشابه \lr{MMC5883MA}‌ بود؛ دلیل تفاوت در مدل انتخاب شده، آدرس متفاوت در پروتکل \lr{I2C} بود تا قرار دادن تمام حسگر ها در یک خط داده، حجم مسیر کشی برد مدار چاپی کاهش یابد.
در طراحی ارائه شده،این حسگر در مرکز آهنربای رینگی و در راستای محور عمودی قرار داده شد. این موقعیت باعث می‌شد حسگر به تغییرات میدان مغناطیسی ناشی از چرخش یا گشتاور حول محور عمودی حساسیت بیشتری داشته باشد. بنابراین، ایده‌ی این جاگذاری، آشکارسازی لنگر انداختن جسم یا گشتاورهای پیچشی اعمال شده بر سطح کشسان بود. هر چند به دلیل نداشتن حسگر مرجع برای گشتاور و عدم کالیبراسیون حسگر برای این مولفه، این ادعا در پژوهش حاضر اثبات نشده و به کارهای پیشرو موکول شد.
	\item
	\textbf{میکروفون  \lr{ISC43434}:}
این حسگر با دارا بودن مکانیزم پیزوالکتریک برای ثبت ارتعاشات صوتی، نماینده‌ی مناسبی از این روش تبدیل در کاربرد لامسه می‌باشد. در پژوهش حاضر با محاط کردن این حسگر در داخل محفظه‌ی کشسان، عملاً دیگر نمی‌توان آن را میکروفون نامید چرا که به جای حساسیت به صدا‌های محیط، ارتعاشات کوچک سطح تماس که ناشی از پدیده‌هایی مانند لغزش هستند را آشکار می‌سازد.
\end{itemize}

یکپارچه‌سازی مکانیزم‌های حسی با پهنای باندهای کاملاً متفاوت (از داده‌های استاتیک تا سیگنال صوتی با نرخ نمونه‌برداری 16 کیلوهرتز)، چالش جدی جمع‌آوری، همگام‌سازی و انتقال داده‌ها را به همراه داشت. برای حل این مشکل، یک برد جمع‌آوری داده اختصاصی مبتنی بر میکروکنترلر \lr{STM32F405RGT6} طراحی و ساخته شد. نوآوری کلیدی و برجسته این برد، پیاده‌سازی یک پروتکل ارتباطی ترکیبی \lr{USB} بود. این راهکار به سیستم‌عامل اجازه می‌دهد تا برد را به طور همزمان به عنوان دو دستگاه مجزا شناسایی کند: یک دستگاه صوتی برای انتقال بی‌وقفه و استاندارد داده‌های ارتعاشی سطح تماس در قالب صدا، و یک پورت سریال مجازی برای انتقال داده‌های استاتیک و کندتر سایر حسگرها. این معماری، ضمن ساده‌سازی فرآیند ارتباط با کامپیوتر، انتقال داده را به بهینه‌ترین شکل ممکن مدیریت کرد.

\subsection{فرآیند کالیبراسیون و مدل‌سازی حسگر لامسه}

فرآیند کالیبراسیون حسگر چندوجهی به دلیل ماهیت غیرخطی و اندرکنش پیچیده بین مبدل‌ها، با یک رویکرد دو مرحله‌ای هوشمند و نوآورانه انجام شد. ر مرحله اول، یک مدل شبکه عصبی تمام‌متصل\LTRfootnote{fully-connected} با استفاده از داده‌های استاتیک و برچسب‌های گسسته آموزش داده شد تا تخمین اولیه‌ای از نیرو به صورت استاتیکی ارائه دهد. سپس خروجی نویزدار این مدل استاتیک، به عنوان ورودی یک فیلتر کالمن مورد استفاده قرار گرفت. فیلتر کالمن صرفاً برای حذف نویز استفاده شد و یک سری زمانی پیوسته، صاف و با کیفیت بالا از نیرو تولید کرد. این خروجی، به عنوان برچسب‌های دینامیکی برای مرحله بعد عمل نمود. در نهایت، این برچسب‌های دینامیکی باکیفیت برای آموزش یک مدل نهایی و قدرتمندتر با معماری حافظه بلند کوتاه مدت \LTRfootnote{Long Short-term Memory} به کار گرفته شدند. این مدل قادر بود به طور مستقیم از داده‌های خام حسگر، نویز را حذف کرده و سیستم را با خورجی کم نویزی مدل‌سازی کند. با این رویکرد ترکیبی، به دقت نهایی ۱۲ میلی‌نیوتن برای نیروی عمودی و ۲۵ میلی‌نیوتن برای نیروی مماسی دست یافتیم که نشان‌دهنده بهبود قابل توجه در پایداری و دقت تخمین نیرو بود.

\subsection{پیاده‌سازی و ارزیابی عملی}
در گام نهایی، عملکرد حسگر چندوجهی کالیبره شده در یک سناریوی واقعی، بر روی یک چنگک رباتیک دو انگشتی ارزیابی شد. در این آزمایش، حسگر با موفقیت توانست نیروهای اعمالی برای گرفتن یک جسم نرم را به صورت بی‌درنگ اندازه‌گیری کرده و مهم‌تر از آن، با تحلیل همزمان داده‌های ارتعاشی ثبت شده توسط میکروفون داخلی، لحظه شروع لغزش را با دقت و سرعت بالایی تشخیص دهد. این دستاورد عملی، موفقیت‌آمیز بودن کل مسیر پژوهش، از ایده اولیه تا پیاده‌سازی نهایی را به اثبات رساند.

\section{چالش ها}

مسیر توسعه و ساخت حسگر چندوجهی، با چالش‌های فنی و عملی متعددی همراه بود. غلبه بر این موانع علاوه بر دستیابی به اهداف نهایی پروژه، منجر به کسب تجربیات ارزشمند و توسعه راهکارهای مهندسی خلاقانه‌ای شد که در ادامه به تفصیل شرح داده می‌شوند.

\subsection{بهبود پاسخ زمانی و پسماند حسگر بارومتریک}
در آزمایش‌های اولیه با حسگر بارومتریک \lr{MPL3115A2}، دو مشکل اساسی پاسخ زمانی کند و پسماند بالا مشاهده شد. تحلیل‌ها نشان داد که این مشکلات ناشی از کوچک بودن منفذ هوا روی محفظه فلزی حسگر است که انتقال فشار از لایه سیلیکونی به دیافراگم داخلی را با تأخیر و به صورت غیربهینه انجام می‌دهد. برای حل این چالش، یک فرآیند اصلاح مکانیکی دقیق و حساس انجام شد. در این فرآیند، ابتدا فضای داخلی حسگر به کمک دیک خلأ با سیلیکون پر شد و پس از خشک شدن، محفظه فلزی به صورت کنترل‌شده تا دمای ۱۲۰ درجه سانتی‌گراد حرارت داده شد تا چسب نگه‌دارنده آن رها شود. سپس محفظه با دقت از روی برد جدا، سوراخ هوای آن با استفاده از مته ۲.۵ میلی‌متری گشادتر شد و مجدداً در جای خود نصب گردید. تزریق سیلیکون پیش از این فرایند یک اقدام احتیاطی هوشمندانه با توجه به حساسیت بالای مدارات داخلی حسگر به تنش مکانیکی بوده. این اصلاحات منجر به بهبود چشمگیر کوپلینگ مکانیکی، افزایش حساسیت و حل کامل مشکلات پاسخ زمانی و برگشت‌پذیری حسگر شد که نتایج کامل آن در فصل \ref{chap2} ارائه گردید. درنهایت این حسگر لامسه‌ی تکی به دقت $2.8mN$ و خطای حدوداً $1\%$ دست یافت؛ مقادیر دقیق صحت در یازه‌های مختلف در جدول \ref{tab:baro_acc} گزارش شده است. 

\subsection{آسیب‌پذیری آرایه سه‌بعدی حسگرهای بارومتریک} 
تلاش برای ساخت یک آرایه سه‌تایی از حسگرهای بارومتریک برای اندازه‌گیری نیروی سه‌محوره، با چالش تکرارپذیری فرآیند اصلاح مکانیکی مواجه شد. به دلیل ماهیت بسیار ظریف و حساس این عملیات، بسیاری از مبدل‌ها در طی فرآیند برداشتن و اصلاح محفظه فلزی آسیب دیده و دچار اشباع دائمی می‌شدند. پس از تلاش‌های مکرر و ناموفق، تصمیم گرفته شد که این رویکرد به طور موقت کنار گذاشته شود تا در کارهای پیش رو، با استفاده از مدل‌های دیگری از حسگرهای بارومتریک که نیاز به چنین اصلاحات تهاجمی ندارند، مجدداً پیگیری شود.

\subsection{ارسال همزمان و یکپارچه‌ی اطلاعات استاتیک و ارتعاش سطح تماس}
یکی از بزرگترین چالش‌های نرم‌افزاری پروژه، پیاده‌سازی یک پروتکل ارتباطی \lr{USB} ترکیبی برای انتقال همزمان داده‌های ارتعاش سطح تماس با ماهیت صوتی (از میکروفون) و سریال (از سایر حسگرها) بود. پکیج‌های نرم‌افزاری و کتابخانه‌های استاندارد موجود، مانند \lr{HAL} شرکت \lr{ST} یا کتابخانه‌های متداول مانند \lr{TinyUSB}، برای این کاربرد خاص مناسب نبودند. اغلب این کتابخانه‌ها برای پیاده‌سازی یک کلاس \lr{USB} واحد (مثلاً فقط سریال یا فقط صوتی) طراحی شده‌اند و معماری لازم برای مدیریت همزمان چند کلاس متفاوت را ندارند. از سوی دیگر، کتابخانه‌هایی مانند \lr{TinyUSB} عمدتاً برای کنترلرهای \lr{USB} کلاسیک طراحی شده‌اند و با معماری پیچیده‌تر
 \LR{USB OTG }
 \LTRfootnote{On-The-Go} 
 که در میکروکنترلرهای سری \lr{STM32F4} به کار رفته است، سازگاری کامل ندارند. در نهایت، برای غلبه بر این محدودیت‌ها، لایه نرم‌افزاری که شامل توصیف‌گرها
 \LTRfootnote{Descriptors}
 و توابع فراخوان
 \LTRfootnote{Callback functions}
 بوده، بر روی میکروکنترلر از پایه نوشته شد. این راهکار، اگرچه زمان‌بر بود، اما انعطاف‌پذیری کامل را برای تعریف یک دستگاه ترکیبی با قابلیت ارسال همزمان داده بر روی کانال‌های منطقی مجزا فراهم آورد.

\subsection{اشباع حسگرهای اثر هال}
چالشی دیگر در مرحله طراحی سخت‌افزار حسگر چندوجهی، عدم تطابق بازه اندازه‌گیری حسگرهای اثر هال با شدت میدان مغناطیسی آهنربای انتخابی بود. به دلیل تنوع بسیار محدود حسگرهای اثر هال در بازار داخلی، مدل‌های انتخابی دارای بازه اندازه‌گیری به مراتب کمتری نسبت به قدرت میدان آهنربا‌های نئودیمیومی بودند. این امر موجب می‌شد که حسگرها بلافاصله پس از قرارگیری آهنربا در نزدیکی آن‌ها به حالت اشباع بروند و قادر به ثبت هیچ تغییری نباشند. برای حل این مشکل، از روش تضعیف حرارتی کنترل‌شده استفاده شد. اغلب مدل‌های آهنربا به خصوص انواع نئودنیومی، در دما‌های بالا شدت شار مغناطیسیشان کاهش می‌یابد\cite{osman2016temperature}. در همین راستا، آهنربا به صورت مرحله به مرحله و با دقت تحت حرارت قرار گرفت تا از طریق پدیده تضعیف مغناطیسی، شدت میدان آن به اندازه‌ای کاهش یابد که در بازه کاری خطی حسگرهای اثر هال قرار گیرد. البته این راهکار، حساسیت سیستم به اغتشاشات محیطی و میدان مغناطیسی زمین را افزایش داد. این مشکل ثانویه نیز با پیاده‌سازی یک روال صفر کردن اولیه\LTRfootnote{Tare} در نرم‌افزار حل شد که در ابتدای هر بار روشن شدن، میدان مغناطیسی محیط را به عنوان نقطه مبنا اندازه‌گیری و از محاسبات بعدی حذف می‌کرد.
\section{نوآوری ها}
این پژوهش با هدف فراتر رفتن از محدودیت‌های حسگرهای لامسه تک‌وجهی، مجموعه‌ای از نوآوری‌های سخت‌افزاری و نرم‌افزاری را ارائه داده‌است که در نهایت به طراحی و ساخت یک حسگر چندوجهی یکپارچه، هوشمند و زیست‌الهام منجر گردیده. نوآوری‌های اصلی این تحقیق بر دو محور کلیدی استوار است: پیاده‌سازی یک رویکرد چندوجهی نوین و ارائه یک راهکار یکپارچه برای سنجش و انتقال داده‌های لامسه.
\subsection{پیاده‌سازی و ارزیابی رویکرد چندوجهی ترکیبی}

یکی از نوآوری‌های بنیادین این پژوهش، ارائه و پیاده‌سازی عملی یک رویکرد ترکیبی است که نشان می‌دهد ادغام هوشمندانه مکانیزم‌های حسی متداول می‌تواند به ساخت یک حسگر لامسه با قابلیت‌های گسترده‌تر منجر شود. در حالی که پژوهش‌های پیشین عمدتاً بر بهینه‌سازی یک نوع مکانیزم (مثلاً پیزومقاومتی یا خازنی) متمرکز بوده‌اند، این تحقیق به بررسی پتانسیل حاصل از ادغام سه مکانیزم متفاوت پرداخت:
\begin{itemize}
	\item \textbf{بارومتریک: }
برای اندازه‌گیری دقیق نیروی استاتیک عمودی.
	\item \textbf{اثر هال: }
برای تفکیک نیروهای برشی و گشتاورهای مماسی.
	\item \textbf{پیزوالکتریک(میکروفون): }
برای ثبت اطلاعات ارتعاشی سطح تماس که ناشی از لغزش و یا پدیده های مشابه هستند.
\end{itemize}

خروجی این رویکرد، دستیابی به یک درک جامع و غنی از تعامل فیزیکی، مشابه حس لامسه انسان یوده‌است. این سیستم ترکیبی، ضعف‌های هر مکانیزم را با نقاط قوت دیگری پوشش داده و یک سیستم حسگر لامسه‌ی مقاوم  و قابل اعتماد را ایجاد کرد.
\subsection{نوآوری در سنجش و انتقال یکپارچه ارتعاشات لامسه}

مهم‌ترین و برجسته‌ترین نوآوری این پژوهش، ارائه یک راهکار سخت‌افزاری-نرم‌افزاری یکپارچه برای سنجش و انتقال اطلاعات استاتیکی و فرکانس بالای تماس، مانند لغزش است. این نوآوری از دو بخش مکمل تشکیل شده است که در کنار یکدیگر، یک سیستم کامل را شکل داده‌اند. در این پژوهش، برای اولین بار یک میکروفون ممز به صورت یکپارچه در داخل ماده کشسان (سیلیکون) همراه با سایر حسگر‌ها و مکانیزم ها تعبیه شد تا به عنوان یک حسگر ارتعاشات لامسه عمل کند. این رویکرد، یک جهش کیفی نسبت به روش‌های پیشین بوده که در آن‌ها داده‌هایی با ماهیت متفاوت، از منابع جدا و منفک دریافت می‌شدند
\cite{yu2023mimictouch}.
همچنین، قرارگیری میکروفون در دل ساختار نرم و کشسان، مانند یک عایق آکوستیک طبیعی عمل کرده و آن را از صداها و ارتعاشات محیطی ایزوله کرد. این امر باعث می‌شود که سیگنال ثبت‌شده تقریباً به طور انحصاری مربوط به ارتعاشات ساختاری ناشی از سطح تماس باشد که منجر به نسبت سیگنال به نویز (SNR) بسیار بالایی می‌شود. همچنین، برخلاف میکروفون‌های خارجی که ارتعاشات را با واسطه هوا و با تضعیف زیاد دریافت می‌کنند، در این طراحی، میکروفون به صورت مکانیکی و مستقیم با  سطح تماس حسگر با اشیاء، کوپل شده است. این اتصال مستقیم، حساسیت سیستم به کوچک‌ترین ارتعاشات ناشی از لغزش‌های اولیه را به شدت افزایش داد.
\\
 یکپارچه سازی حسگر میکروفون به عنوان گیرنده‌ی ارتعاشات فرکانس بالا در کنار حسگر‌های بارومتریک و اثر هال، چالش جدیدی را به وجود آورد که آن انتقال همزمان و پایدار جریان داده‌های ارتعاش با نرخ نمونه‌برداری 16 کیلوهرتز در کنار داده‌های استاتیک حسگرهای بارومتریک و اثر هال با نرخ 100 هرتز بود. برای حل این چالش، یک برد جمع‌آوری داده اختصاصی با یک پروتکل ارتباطی منحصربه‌فرد طراحی شد. این برد با استفاده از پروتکل USB و کلاس ترکیبی در میکروکنترلر STM32، خود را به طور همزمان به عنوان دو دستگاه مجزا به رایانه معرفی می‌کند:
\begin{itemize}
	\item 
	دستگاه صوتی استاندارد برای انتقال بی‌وقفه و پرسرعت داده‌های میکروفون از طریق درگاه استاندارد صوتی سیستم‌عامل.
	\item 
	پورت سریال مجازی برای انتقال داده‌های مربوط به حسگرهای بارومتریک و اثر هال.
\end{itemize}

\begin{figure}[t]
	\centering
	\centerline{\includegraphics[width=0.8\textwidth]{barotac_triple}}
	\caption[برد جمع‌آوری داده در \LR{Device manager}]{
		برد جمع‌آوری داده در
		 \LR{Device manager}‌.
	}
	\label{fig:DAQ_devicemanager}
\end{figure}

این معماری نوآورانه، نیاز به همگام‌سازی پیچیده نرم‌افزاری در سمت کامپیوتر را از بین برده و امکان انتقال پایدار و بهینه‌ی داده‌های چندوجهی را تنها از طریق یک کابل USB واحد و بدون نیاز به هیچ‌گونه مبدل یا درایور خارجی فراهم ‌آورد. در واقع، این دو نوآوری سخت‌افزاری و نرم‌افزاری در کنار هم، یک راهکار کامل برای ادغام وجوه حسی متنوع به ربات‌ها به شیوه‌ای کارآمد، ارزان و قابل اعتماد فراهم کرد.
\section{کارهای  پیش‌رو}
دستاوردهای پژوهش حاضر، ضمن ارائه یک راهکار نوین در زمینه حسگری لامسه، زمینه‌ساز مسیرهای تحقیقاتی آتی و فرصت‌های جدید برای توسعه در حوزه رباتیک تعاملی گردیده است. این تحقیق نقطه پایانی بر این مسیر نبوده و خود سرآغازی برای افق‌های گسترده‌تر است. در ادامه، مجموعه‌ای از مسیرهای تحقیقاتی و توسعه‌ای که به عنوان گام‌های بعدی این پروژه تعریف شده‌اند و برخی از آن‌ها در زمان نوشتن این گزارش نیز در حال پیگیری هستند، معرفی و شرح داده شده است.
\subsection{بهینه‌سازی طراحی با بهره‌گیری از قطعات الکترونیکی جدید تر}

یکی از نتایج کلیدی حاصل از این پژوهش، شناسایی تأثیر مستقیم مشخصات فنی قطعات بر پیچیدگی فرآیند ساخت و عملکرد نهایی حسگر بود. بسیاری از چالش‌های مهندسی، مانند نیاز به اصلاحات مکانیکی بر روی حسگر بارومتریک یا پدیده اشباع در حسگرهای اثر هال، مستقیماً به محدودیت‌های فناوری قطعات موجود در زمان طراحی بازمی‌گشت. پیشرفت‌های اخیر در فناوری نیمه‌هادی‌ها، اکنون قطعاتی با قابلیت‌های برتر را معرفی کرده‌اند که می‌توانند این چالش‌ها را مرتفع سازند. به طور مثال، حسگرهای بارومتریک نظیر \lr{LPS22HHTR} از شرکت \lr{STMicroelectronics} (شکل \ref{fig:lps22hhtr})دارای ساختاری هستند که در آن دیافراگم حسگر با یک محافظ سیلیکونی مستقیماً با محیط در ارتباط است. این طراحی، نیاز به هرگونه اصلاح مکانیکی را حذف کرده، ریسک آسیب به قطعه را به حداقل می‌رساند و فرآیند ساخت را به میزان قابل توجهی ساده‌تر و تکرارپذیرتر  می‌نماید.
\begin{figure}[t]
	\centering
	\centerline{\includegraphics[width=0.6\textwidth]{lps22hhtr}}
	\caption[حسگر بارومتر \lr{LPS22HHTR}.]{
		حسگر بارومتر
		 \lr{LPS22HHTR}
		 .
	}
	\label{fig:lps22hhtr}
\end{figure}

همچنین حسگرهای اثر هال دیجیتال با بازه‌های اندازه‌گیری وسیع‌تر که پس از طراحی ارائه شده، در دسترس قرار گرفته اند، چالش اشباع شدن در حضور میدان‌های مغناطیسی قوی را برطرف نموده و نیاز به فرآیند پیچیده و غیرقایل اطمینان تضعیف کنترل‌شده میدان از طریق حرارت را از بین می‌برند. این امر به طور مستقیم به بهبود صحت و افزایش دقت در اندازه‌گیری نیروهای برشی منجر خواهد شد.
\begin{figure}[t]
	\centering
	\centerline{\includegraphics[width=0.6\textwidth]{TMAG5273A2}}
	\caption[حسگر اثر هال \lr{TMAG5273A2}.]{
		حسگر اثر هال
		\lr{TMAG5273A2}
		.
	}
	\label{fig:tmag527a2}
\end{figure}
بازطراحی ساختار حسگر بر پایه قطعات جدیدتر، ضمن حفظ اصول عملکردی فعلی، به تولید یک نسخه با استحکام، پایایی و قابلیت تولید در مقیاس صنعتی منجر خواهد شد
\subsection{توسعه زیرساخت کالیبراسیون جامع با مرجع نیروی شش‌محوره}
زیرساخت کالیبراسیون مورد استفاده در این پژوهش امکان تفکیک نیروهای نرمال و مماسی در یک راستا را فراهم آورد و برای کاربر‌های برداشت اشیاء کافی بود. اما جهت بهره‌برداری از ظرفیت کامل یک حسگر چندوجهی، استفاده از یک سیستم کالیبراسیون جامع‌تر امری ضروری است. در تعاملات فیزیکی واقعی، نیروها و گشتاورها به صورت ترکیبی و تزویج‌شده بر سطح حسگر اعمال می‌شوند. به منظور مدل‌سازی دقیق این اندرکنش‌های پیچیده و شناسایی اثرات متقابل بین محورهای مختلف، استفاده از یک حسگر نیرو-گشتاور شش‌محوره به عنوان مرجع کالیبراسیون اجتناب‌ناپذیر است.
طراحی یک مجموعه کالیبراسیون جدید که قادر به اعمال دقیق بردارها و گشتاورهای کنترل‌شده در فضای سه‌بعدی باشد، امکان استخراج مدل کالیبراسیون کامل و غیرخطی بسیار دقیق‌تری را فراهم می‌آورد. این امر حسگر را برای تخمین دقیق بردار کامل نیرو و گشتاور توانمند ساخته و آن را برای کاربردهای پیچیده دستکاری ماهرانه
\LTRfootnote{Dexterous Manipulation}
آماده می‌سازد.
\subsection{پیاده‌سازی الگوریتم‌های پیشرفته پردازش سیگنال و همجوشی داده}

سیگنال ارتعاشی حاصل از میکروفون ممز، منبعی غنی از اطلاعات دینامیکی با پهنای باند بالا است که در این پژوهش صرفاً به صورت مقدماتی مورد تحلیل قرار گرفت. تحقیقات آتی می‌توانند بر روی استخراج ویژگی‌های عمیق‌تر از این سیگنال با استفاده از تکنیک‌های پیشرفته متمرکز شوند. به طور مثال، به کار گیری روش های تحلیل سیگنال در حوزه زمان-فرکانس مانند تبدیل فوریه زمان-کوتاه
\LTRfootnote{Short time Fourier transfrom (STFT)}
یا تحلیل موجک پیوسته
\LTRfootnote{Continious Wavelet transform (CWT)}
، امکان شناسایی الگوهای فرکانسی گذرا که مشخصه رویدادهای دینامیکی هستند را فراهم می‌آورد. در پژوهش های
\cite{chang2024investigation,yu2023mimictouch}
 حتی بافت سطوح مختلف نیز بر اساس توزیع انرژی سیگنال در باندهای فرکانسی متفاوت، طبقه‌بندی شده است.

در این پژوهش، همجوشی در سطح داده خام پیاده‌سازی شد. مسیرهای آتی می‌توانند به بررسی معماری‌های همجوشی در سطح ویژگی و سطح تصمیم  بپردازند. برای نمونه، می‌توان ویژگی‌های استخراج‌شده از داده‌های استاتیک (مانند مؤلفه‌های بردار نیرو) را با ویژگی‌های حوزه فرکانس ادغام کرده و به یک مدل یادگیری ماشین جهت طبقه‌بندی دقیق وضعیت تماس (مانند گرفتن پایدار، لغزش، غلتش) ارائه داد.

\subsection{تلفیق با مودالیته‌های حسگری مکمل جهت کاربردهای تعاملی پیشرفته}

حسگر لامسه طراحی‌شده می‌تواند به عنوان مؤلفه‌ای کلیدی در یک سیستم ادراکی گسترده‌تر عمل کند. با افزودن مودالیته‌های دیگر، می‌توان کاربرد آن را به حوزه‌های نوظهوری نظیر یادگیری تقلیدی
 \LTRfootnote{Imitation Learning}
 گسترش داد. به طور مثال با تعبیه حسگر بر روی یک دستکش ابزار دقیق مجهز به واحدهای اندازه‌گیری اینرسی
 \LTRfootnote{IMUs} 
 یا به‌کارگیری سیستم‌های ردیابی حرکت
  \LTRfootnote{Motion Capture}
  ، می‌توان داده‌های لامسه را با اطلاعات سینماتیکی دست انسان (موقعیت، جهت‌گیری و سرعت مفاصل) به صورت همزمان ثبت کرد. این مجموعه داده غنی و چندوجهی به ربات امکان می‌دهد تا سیاست‌های پیچیده دستکاری را بیاموزد؛ سیاست‌هایی که در آن‌ها نه تنها مقدار نیروی اعمالی، بلکه نحوه توزیع زمانی و مکانی آن در طول اجرای یک وظیفه نیز حائز اهمیت است.
\subsection{ طراحی آرایه‌های حسگری ماژولار و شبکه‌ای}
چشم‌انداز غایی این فناوری، گذار از حسگرهای منفرد به سمت یک سطح حسگر توزیع‌شده و گسترده، موسوم به پوست رباتیک است. حسگر چندوجهی این پژوهش به دلیل ماهیت یکپارچه خود، یک واحد سازنده ایده‌آل برای تحقق این چشم‌انداز می‌باشد. کوچک‌سازی واحد حسگر و پیاده‌سازی آن بر روی بردهای مدار چاپی انعطاف‌پذیر یکی از اقدامات در این راستا می‌باشد. این امر امکان ساخت آرایه‌های متراکم و انعطاف‌پذیری را فراهم می‌کند که قادرند با سطوح منحنی و پیچیده هندسی (مانند انگشتان یا کف دست ربات) تطبیق یابند. چنین آرایه‌ای، یک نقشه توزیع‌شده از فشار و برش با وضوح فضایی بالا را برای ربات فراهم می‌آورد.

مدیریت حجم عظیم داده‌های تولید شده توسط صدها واحد حسگری، یک چالش اساسی است. راهکار آینده، شبکه‌سازی ماژول‌های حسگری بر اساس پروتکل هایی مانند 
\LR{CAN bus}
است که قابلیت اتصال و انتقال داده از چندین دستگاه را فراهم می‌کنند. علاوه بر این، جهت جلوگیری از ایجاد بار محاسباتی سنگین، می‌توان از یک معماری پردازش سلسله‌مراتبی
 \LTRfootnote{Hierarchical Processing} 
 بهره برد. در این معماری، هر واحد سازنده دارای پردازش محلی است که وظیفه پیش‌پردازش و استخراج ویژگی‌های اولیه را بر عهده گرفته و تنها اطلاعات کلیدی و فشرده‌شده را به لایه‌های بالاتر پردازشی ارسال می‌کند. این رویکرد، ضمن کاهش بار محاسباتی، مقیاس‌‌پذیری سیستم را به طور قابل توجهی افزایش می‌‌دهد.
